%
% this file is encoded in utf-8
% v2.0 (Apr. 5, 2009)
% do not change the content of this file
% unless the thesis layout rule is changed
% 無須修改本檔內容，除非校方修改了
% 封面、書名頁、中文摘要、英文摘要、誌謝、目錄、表目錄、圖目錄、符號說明
% 等頁之格式

% make the line spacing in effect
\renewcommand{\baselinestretch}{\mybaselinestretch}
\large % it needs a font size changing command to be effective

% default variables definitions
% 此處只是預設值，不需更改此處
% 請更改 my_names.tex 內容
\newcommand\cTitle{論文題目}
\newcommand\eTitle{MY THESIS TITLE}
\newcommand\myCname{王鐵雄}
\newcommand\myEname{Aron Wang}
\newcommand\myStudentID{M9315048}    % add student id by saiba 
\newcommand\advisorCnameA{南宮明博士}
\newcommand\advisorEnameA{Dr.~Ming Nangong}
\newcommand\advisorCnameB{李斯坦博士}
\newcommand\advisorEnameB{Dr.~Stein Lee}
\newcommand\advisorCnameC{徐　石博士}
\newcommand\advisorEnameC{Dr.~Sean~Hsu}
\newcommand\univCname{元智大學}
\newcommand\univEname{Yuan Ze University}
\newcommand\deptCname{光電工程研究所}
\newcommand\fulldeptEname{Graduate School of Electro-Optical Engineering}
\newcommand\deptEname{Electro-optical Engineering}
\newcommand\collEname{College of Engineering}
\newcommand\degreeCname{碩士}
\newcommand\degreeEname{Master of Science}
\newcommand\cYear{九十四}
\newcommand\cMonth{六}
\newcommand\cDay{六}
\newcommand\eYear{2006}
\newcommand\eMonth{June}
\newcommand\ePlace{Chungli, Taoyuan, Taiwan}


 % user's names; to replace those default variable definitions
%
% this file is encoded in utf-8
% v2.0 (Apr. 5, 2009)
% 填入你的論文題目、姓名等資料
% 如果題目內有必須以數學模式表示的符號，請用 \mbox{} 包住數學模式，如下範例
% 如果中文名字是單名，與姓氏之間建議以全形空白填入，如下範例
% 英文名字中的稱謂，如 Prof. 以及 Dr.，其句點之後請以不斷行空白~代替一般空白，如下範例
% 如果你的指導教授沒有如預設的三位這麼多，則請把相對應的多餘教授的中文、英文名
%    的定義以空的大括號表示
%    如，\renewcommand\advisorCnameB{}
%          \renewcommand\advisorEnameB{}
%          \renewcommand\advisorCnameC{}
%          \renewcommand\advisorEnameC{}

% 論文題目 (中文)
\renewcommand\cTitle{%
3D 點雲與SIFT特徵點室內影像定位之研究
}

% 論文題目 (英文)
\renewcommand\eTitle{%
3D Point Cloud and SIFT Descriptor Indoor Localization Research
}

% 我的姓名 (中文)
\renewcommand\myCname{陳致良}

% 我的姓名 (英文)
\renewcommand\myEname{ZZhi-Liang Chen}

% 我的學號 (英文)
\renewcommand\myStudentID{M9915057}

% 指導教授A的姓名 (中文)
\renewcommand\advisorCnameA{項天瑞博士}

% 指導教授A的姓名 (英文)
\renewcommand\advisorEnameA{Dr. Tien-Ruey Hsiang}

% 指導教授B的姓名 (中文)
\renewcommand\advisorCnameB{}

% 指導教授B的姓名 (英文)
\renewcommand\advisorEnameB{}

% 指導教授C的姓名 (中文)
\renewcommand\advisorCnameC{}

% 指導教授C的姓名 (英文)
\renewcommand\advisorEnameC{}

% 校名 (中文)
\renewcommand\univCname{國立台灣科技大學}

% 校名 (英文)
\renewcommand\univEname{National Taiwan University of Science and Technology}

% 系所名 (中文)
\renewcommand\deptCname{資訊工程研究所}

% 系所全名 (英文)
\renewcommand\fulldeptEname{Department of Computer Science and Information Engineering}

% 系所短名 (英文, 用於書名頁學位名領域)
%\renewcommand\deptEname{Electro-Optical Engineering}

% 學院英文名 (如無，則以空的大括號表示)
\renewcommand\collEname{College of Electrical Engineering and Computer Science}

% 學位名 (中文)
\renewcommand\degreeCname{碩士}

% 學位名 (英文)
\renewcommand\degreeEname{Master of Science}

% 口試年份 (中文、民國)
\renewcommand\cYear{一百零一}

% 口試月份 (中文)
\renewcommand\cMonth{七}

% 口試日期 (中文)
\renewcommand\cDay{十}

% 口試年份 (阿拉伯數字、西元)
\renewcommand\eYear{2012}

% 口試月份 (英文)
\renewcommand\eMonth{July}

% 學校所在地 (英文)
%\renewcommand\ePlace{Taipei City, Taiwan}

%畢業級別；用於書背列印；若無此需要可忽略
\newcommand\GraduationClass{98}

%%%%%%%%%%%%%%%%%%%%%%

% 使用 hyperref 在 pdf 簡介欄裡填入相關資料
\ifx\hypersetup\undefined
	\relax  % do nothing
\else
	\hypersetup{
	pdftitle=\cTitle,
	pdfauthor=\myCname}
\fi
	

\newcommand\itsempty{}
%%%%%%%%%%%%%%%%%%%%%%%%%%%%%%%
%       NTUST cover 封面
%%%%%%%%%%%%%%%%%%%%%%%%%%%%%%%
%
\begin{titlepage}
% no page number
% next page will be page 1

% aligned to the center of the page
\begin{center}
% font size (relative to 12 pt):
% \large (14pt) < \Large (18pt) < \LARGE (20pt) < \huge (24pt)< \Huge (24 pt)
%

\begin{figure}[htbp]
	\begin{minipage}[b]{5cm} 
		\raggedright
		\includegraphics[width=1.5in]{frontpages/ntust_logo}
		\label{fig:ntust_logo}
	\end{minipage}% 
	\begin{minipage}[b]{0.5\textwidth} 
	\centering
	\makebox[3cm][c]{\Huge{\univCname}}\\  %顯示中文校名
	\vspace{0.5cm}
	\makebox[3cm][c]{\Huge{\deptCname}}\\ %顯示中文系所名
	\vspace{0.5cm}
	\end{minipage}%
\\ 
\rule{16cm}{3pt}
\end{figure}
%\hfill

\vspace{1cm}
\makebox[6cm][s]{\textbf{\Huge{\degreeCname 論文}}}\\ %顯示論文種類 (中文)
\vspace{1cm}
%
% Set the line spacing to single for the titles (to compress the lines)
\renewcommand{\baselinestretch}{1}   %行距 1 倍
%\large % it needs a font size changing command to be effective
\Large{\cTitle}\\  % 中文題目
%
\vspace{1cm}
%
\Large{\eTitle}\\ %英文題目
\vspace{5cm}
% \makebox is a text box with specified width;
% option s: stretch
% use \makebox to make sure
% 「研究生：」 與「指導教授：」occupy the same width
%
% remark by saiba
% 由於全型符號在xelatex中無法正確的左右分散，所以將中文與符號部分分開處理

\hspace{4cm} \makebox[3cm][s]{\Large{研究生}}\makebox[.5cm][s]{\Large{：}}
\Large{\myCname}  % 顯示作者中文名
\hfill \makebox[1cm][s]{}\\
%
\vspace{0.3cm}
\hspace{4cm} \makebox[3cm][s]{\Large{學號}}\makebox[.5cm][s]{\Large{：}}
\Large{\myStudentID}  %顯示指導教授A中文名
\hfill \makebox[1cm][s]{}\\
%
\vspace{1cm}
\hspace{4cm} \makebox[3cm][s]{\Large{指導教授}}\makebox[.5cm][s]{\Large{：}}
\Large{\advisorCnameA}  %顯示指導教授A中文名
\hfill \makebox[1cm][s]{}\\
%
% 判斷是否有共同指導的教授 B
\ifx \advisorCnameB  \itsempty
\relax % 沒有 B 教授，所以不佔版面，不印任何空白
\else
% 共同指導的教授 B
\hspace{4.5cm} \makebox[3cm][s]{}
\Large{\advisorCnameB}  %顯示指導教授B中文名
\hfill \makebox[1cm][s]{}\\
\fi
%
% 判斷是否有共同指導的教授 C
\ifx \advisorCnameC  \itsempty
\relax % 沒有 C 教授，所以不佔版面，不印任何空白
\else
% 共同指導的教授 C
\hspace{4.5cm} \makebox[3cm][s]{}
\Large{\advisorCnameC}  %顯示指導教授C中文名
\hfill \makebox[1cm][s]{}\\
\fi
%
\vfill
\makebox[10cm][s]{\Large{中華民國\cYear 年\cMonth 月\cDay 日}}%
%
\end{center}
% Resume the line spacing to the desired setting
\renewcommand{\baselinestretch}{\mybaselinestretch}   %恢復原設定
% it needs a font size changing command to be effective
% restore the font size to normal
\normalsize
\end{titlepage}
%%%%%%%%%%%%%%

%% 從摘要到本文之前的部份以小寫羅馬數字印頁碼
% 但是從「書名頁」(但不印頁碼) 就開始計算
\setcounter{page}{1}
\pagenumbering{roman} 
%\pagenumbering{arabic}



%%%%%%%%%%%%%%%%%%%%%%%%%%%%%%%%
%%       書名頁 
%%%%%%%%%%%%%%%%%%%%%%%%%%%%%%%%
%%
%\newpage
%
%% 判斷是否要浮水印？
%\ifx\mywatermark\undefined 
%  \thispagestyle{empty}  % 無頁碼、無 header (無浮水印)
%\else
%  \thispagestyle{EmptyWaterMarkPage} % 無頁碼、有浮水印
%\fi
%
%%no page number
%% create an entry in table of contents for 書名頁
%\phantomsection % for hyperref to register this
%\addcontentsline{toc}{chapter}{\nameInnerCover}
%
%
%% aligned to the center of the page
%\begin{center}
%% font size (relative to 12 pt):
%% \large (14pt) < \Large (18pt) < \LARGE (20pt) < \huge (24pt)< \Huge (24 pt)
%% Set the line spacing to single for the titles (to compress the lines)
%\renewcommand{\baselinestretch}{1}   %行距 1 倍
%% it needs a font size changing command to be effective
%%中文題目
%\Large{\cTitle}\\ %%%%%
%\vspace{1cm}
%% 英文題目
%\Large{\eTitle}\\ %%%%%
%%\vspace{1cm}
%\vfill
%% \makebox is a text box with specified width;
%% option s: stretch
%% use \makebox to make sure
%% 「研究生：」 與「指導教授：」occupy the same width
%\makebox[3cm][s]{\large{研 究 生：}}
%\makebox[3cm][l]{\large{\myCname}} %%%%%
%\hfill
%\makebox[2cm][s]{\large{Student: }}
%\makebox[5cm][l]{\large{\myEname}}\\ %%%%%
%%
%%\vspace{1cm}
%%
%\makebox[3cm][s]{\large{指導教授：}}
%\makebox[3cm][l]{\large{\advisorCnameA}} %%%%%
%\hfill
%\makebox[2cm][s]{\large{Advisor: }}
%\makebox[5cm][l]{\large{\advisorEnameA}}\\ %%%%%
%%
%% 判斷是否有共同指導的教授 B
%\ifx \advisorCnameB  \itsempty
%\relax % 沒有 B 教授，所以不佔版面，不印任何空白
%\else
%%共同指導的教授B
%\makebox[3cm][s]{}
%\makebox[3cm][l]{\large{\advisorCnameB}} %%%%%
%\hfill
%\makebox[2cm][s]{}
%\makebox[5cm][l]{\large{\advisorEnameB}}\\ %%%%%
%\fi
%%
%% 判斷是否有共同指導的教授 C
%\ifx \advisorCnameC  \itsempty
%\relax % 沒有 C 教授，所以不佔版面，不印任何空白
%\else
%%共同指導的教授C
%\makebox[3cm][s]{}
%\makebox[3cm][l]{\large{\advisorCnameC}} %%%%%
%\hfill
%\makebox[2cm][s]{}
%\makebox[5cm][l]{\large{\advisorEnameC}}\\ %%%%%
%\fi
%%
%% Resume the line spacing to the desired setting
%\renewcommand{\baselinestretch}{\mybaselinestretch}   %恢復原設定
%\large %it needs a font size changing command to be effective
%%
%\vfill
%\makebox[4cm][s]{\large{\univCname}}\\% 校名
%\makebox[6cm][s]{\large{\deptCname}}\\% 系所名
%\makebox[3cm][s]{\large{\degreeCname 論文}}\\% 學位名
%%
%%\vspace{1cm}
%\vfill
%\large{A Thesis}\\%
%\large{Submitted to }%
%%
%\large{\fulldeptEname}\\%系所全名 (英文)
%%
%%
%\ifx \collEname  \itsempty
%\relax % 沒有學院名 (英文)，所以不佔版面，不印任何空白
%\else
%% 有學院名 (英文)
%\large{\collEname}\\% 學院名 (英文)
%\fi
%%
%\large{\univEname}\\%校名 (英文)
%%
%\large{in Partial Fulfillment of the Requirements}\\
%%
%\large{for the Degree of}\\
%%
%\large{\degreeEname}\\%學位名(英文)
%%
%\large{in}\\
%%
%\large{\deptEname}\\%系所短名(英文；表明學位領域)
%%
%\large{\eMonth\ \eYear}\\%月、年 (英文)
%%
%\large{\ePlace}% 學校所在地 (英文)
%\vfill
%\large{中華民國}%
%\large{\cYear}% %%%%%
%\large{年}%
%\large{\cMonth}% %%%%%
%\large{月}\\
%\end{center}
%% restore the font size to normal
%\normalsize
%\clearpage
%%%%%%%%%%%%%%%%%%%%%%%%%%%%%%%%
%%       論文口試委員審定書 (計頁碼，但不印頁碼) 
%%%%%%%%%%%%%%%%%%%%%%%%%%%%%%%%
%%
%% insert the printed standard form when the thesis is ready to bind
%% 在口試完成後，再將已簽名的審定書放入以便裝訂
%% create an entry in table of contents for 審定書
%% 目前送出空白頁
%\newpage%
%{\thispagestyle{empty}%
%\phantomsection % for hyperref to register this
%\addcontentsline{toc}{chapter}{\nameCommitteeForm}%
%\mbox{}\clearpage}

%%%%%%%%%%%%%%%%%%%%%%%%%%%%%%%
%       指導教授推薦書 (計頁碼，但不印頁碼)
%%%%%%%%%%%%%%%%%%%%%%%%%%%%%%%
%
% insert the printed standard form when the thesis is ready to bind
% 在口試完成後，再將已簽名的推薦書放入以便裝訂
% create an entry in table of contents for 推薦書
% 目前送出空白頁
\newpage%
{\thispagestyle{empty}%
\phantomsection % for hyperref to register this
\addcontentsline{toc}{chapter}{\nameRecommForm}%
\mbox{}\clearpage}


%%%%%%%%%%%%%%%%%%%%%%%%%%%%%%%
%       考試委員審定書 (計頁碼，但不印頁碼) 
%%%%%%%%%%%%%%%%%%%%%%%%%%%%%%%
%
% insert the printed standard form when the thesis is ready to bind
% 在口試完成後，再將已簽名的審定書放入以便裝訂
% create an entry in table of contents for 審定書
% 目前送出空白頁
\newpage%
{\thispagestyle{empty}%
\phantomsection % for hyperref to register this
\addcontentsline{toc}{chapter}{\nameCommitteeForm}%
\mbox{}\clearpage}

%%%%%%%%%%%%%%%%%%%%%%%%%%%%%%%
%       中文摘要 
%%%%%%%%%%%%%%%%%%%%%%%%%%%%%%%
%
\newpage
\thispagestyle{plain}  % 無 header，但在浮水印模式下會有浮水印
% create an entry in table of contents for 中文摘要
\phantomsection % for hyperref to register this
\addcontentsline{toc}{chapter}{\nameCabstract}

% aligned to the center of the page
\begin{center}
% font size (relative to 12 pt):
% \large (14pt) < \Large (18pt) < \LARGE (20pt) < \huge (24pt)< \Huge (24 pt)
% Set the line spacing to single for the names (to compress the lines)
\renewcommand{\baselinestretch}{1}   %行距 1 倍
% it needs a font size changing command to be effective
\large{\cTitle}\\  %中文題目
\vspace{0.83cm}
% \makebox is a text box with specified width;
% option s: stretch
% use \makebox to make sure
% each text field occupies the same width
\makebox[1.5cm][s]{\large{學生：}}
\makebox[3cm][l]{\large{\myCname}} %學生中文姓名
\hfill
%
\makebox[3cm][s]{\large{指導教授：}}
\makebox[3cm][l]{\large{\advisorCnameA}} \\ %教授A中文姓名
%
% 判斷是否有共同指導的教授 B
\ifx \advisorCnameB  \itsempty
\relax % 沒有 B 教授，所以不佔版面，不印任何空白
\else
%共同指導的教授B
\makebox[1.5cm][s]{}
\makebox[3cm][l]{} %%%%%
\hfill
\makebox[3cm][s]{}
\makebox[3cm][l]{\large{\advisorCnameB}}\\ %教授B中文姓名
\fi
%
% 判斷是否有共同指導的教授 C
\ifx \advisorCnameC  \itsempty
\relax % 沒有 C 教授，所以不佔版面，不印任何空白
\else
%共同指導的教授C
\makebox[1.5cm][s]{}
\makebox[3cm][l]{} %%%%%
\hfill
\makebox[3cm][s]{}
\makebox[3cm][l]{\large{\advisorCnameC}}\\ %教授C中文姓名
\fi
%
\vspace{0.42cm}
%
\large{\univCname}\large{\deptCname}\\ %校名系所名
\vspace{0.83cm}
%\vfill
\makebox[2.5cm][s]{\large{摘要}}\\
\end{center}
% Resume the line spacing to the desired setting
\renewcommand{\baselinestretch}{\mybaselinestretch}   %恢復原設定
%it needs a font size changing command to be effective
% restore the font size to normal
\normalsize
%%%%%%%%%%%%%
本篇論文探討如何利用單張彩色影像來重建出三維人臉模型。我們的方法是利用彩色影像來取代灰階值去建立張量模型(tensor model)，在人臉資料庫中，是用典型相關分析(Canonical correlation analysis)是來建立彩色影像與深度資訊的對應關係，一旦建立好屬於各自的對應關係後，在重建人臉的過程中，就只需要單張彩色影像即可利用典型相關分析的對應關係來推算出正確的深度資訊。實驗中，我們的方法可以在不同的光線環境跟人臉角度得到不錯的效果。

%%%%%%%%%%%%%%%%%%%%%%%%%%%%%%%
%       英文摘要 
%%%%%%%%%%%%%%%%%%%%%%%%%%%%%%%
%
\newpage
\thispagestyle{plain}  % 無 header，但在浮水印模式下會有浮水印

% create an entry in table of contents for 英文摘要
\phantomsection % for hyperref to register this
\addcontentsline{toc}{chapter}{\nameEabstract}

% aligned to the center of the page
\begin{center}
% font size:
% \large (14pt) < \Large (18pt) < \LARGE (20pt) < \huge (24pt)< \Huge (24 pt)
% Set the line spacing to single for the names (to compress the lines)
\renewcommand{\baselinestretch}{1}   %行距 1 倍
%\large % it needs a font size changing command to be effective
\large{\eTitle}\\  %英文題目
\vspace{0.83cm}
% \makebox is a text box with specified width;
% option s: stretch
% use \makebox to make sure
% each text field occupies the same width
\makebox[2cm][s]{\large{Student: }}
\makebox[5cm][l]{\large{\myEname}} %學生英文姓名
\hfill
%
\makebox[2cm][s]{\large{Advisor: }}
\makebox[5cm][l]{\large{\advisorEnameA}} \\ %教授A英文姓名
%
% 判斷是否有共同指導的教授 B
\ifx \advisorCnameB  \itsempty
\relax % 沒有 B 教授，所以不佔版面，不印任何空白
\else
%共同指導的教授B
\makebox[2cm][s]{}
\makebox[5cm][l]{} %%%%%
\hfill
\makebox[2cm][s]{}
\makebox[5cm][l]{\large{\advisorEnameB}}\\ %教授B英文姓名
\fi
%
% 判斷是否有共同指導的教授 C
\ifx \advisorCnameC  \itsempty
\relax % 沒有 C 教授，所以不佔版面，不印任何空白
\else
%共同指導的教授C
\makebox[2cm][s]{}
\makebox[5cm][l]{} %%%%%
\hfill
\makebox[2cm][s]{}
\makebox[5cm][l]{\large{\advisorEnameC}}\\ %教授C英文姓名
\fi
%
\vspace{0.42cm}
\large{Submitted to }\large{\fulldeptEname}\\  %英文系所全名
%
\ifx \collEname  \itsempty
\relax % 如果沒有學院名 (英文)，則不佔版面，不印任何空白
\else
% 有學院名 (英文)
\large{\collEname}\\% 學院名 (英文)
\fi
%
\large{\univEname}\\  %英文校名
\vspace{0.83cm}
%\vfill
%
\large{ABSTRACT}\\
%\vspace{0.5cm}
\end{center}
% Resume the line spacing the desired setting
\renewcommand{\baselinestretch}{\mybaselinestretch}   %恢復原設定
%\large %it needs a font size changing command to be effective
% restore the font size to normal
\normalsize
%%%%%%%%%%%%%
This paper develops a tensor-based 3D face reconstruction approach from a single color image. Instead of the grayscale image, we also consider additional color factors in constructing the tensor model. Canonical correlation analysis is applied to establish the relationship between the color image and the depth information in the face database. During the face reconstruction, given a single color face image, the depth estimation is computed from the CCA-based mapping between the tensor models. Experimental results show our approach is better suited under different lighting conditions and poses.

%%%%%%%%%%%%%%%%%%%%%%%%%%%%%%%
%       誌謝 
%%%%%%%%%%%%%%%%%%%%%%%%%%%%%%%
%
% Acknowledgment
\newpage
\chapter*{\protect\makebox[5cm][s]{\nameAckn}} %\makebox{} is fragile; need protect
\phantomsection % for hyperref to register this
\addcontentsline{toc}{chapter}{\nameAckn}
本論文能夠完成，首先要感謝的是指導教授項天瑞老師。
老師嚴謹的治學態度，讓我不但在學術研究上學習到更謹慎的思考，也在日常生活上獲益良多。


感謝實驗室的同學們，實驗室的生活有苦有樂，有你們才讓我能撐得下這三年漫長的時間。
感謝建群、松翰學長給我的指導。感謝益偉、崇峰、承誌、誠儀、慶豪學長以及恩緯、盈樽、嘉駿、世寬學長給我的指導。感謝實驗室的學弟妹們，訓哲、立昂、致良、青緯、庭耀、宗博、雅筑，常常幫了我不少忙。更要感謝我的同梯們，貴彥、薇穎、冠佑，我們同甘共苦，一起奮鬥，特別是在弄計畫的過程中，冠佑幫助了我很多，有你們在我才能走到這一步。
也感謝好朋友志強跟我一起修正文法與用詞。

最後我要把我最深的感謝留給我的家人。謝謝我的爸媽跟兩位姊姊，你們讓我沒有經濟壓力地讀完這這碩士學位，也常常給我很多鼓勵，今天我終於拿到這個學位，終於可以讓你們放下心上的一塊大石頭了。

%%%%%%%%%%%%%%%%%%%%%%%%%%%%%%%
%       目錄 
%%%%%%%%%%%%%%%%%%%%%%%%%%%%%%%
%
% Table of contents
\newpage
\renewcommand{\contentsname}{\protect\makebox[5cm][s]{\nameToc}}
%\makebox{} is fragile; need protect
\phantomsection % for hyperref to register this
\addcontentsline{toc}{chapter}{\nameToc}
\tableofcontents

%%%%%%%%%%%%%%%%%%%%%%%%%%%%%%%
%       圖目錄 
%%%%%%%%%%%%%%%%%%%%%%%%%%%%%%%
%
% List of Figures
\newpage
\renewcommand{\listfigurename}{\protect\makebox[5cm][s]{\nameTof}}
%\makebox{} is fragile; need protect
\phantomsection % for hyperref to register this
\addcontentsline{toc}{chapter}{\nameTof}
\listoffigures


%%%%%%%%%%%%%%%%%%%%%%%%%%%%%%%
%       表目錄 
%%%%%%%%%%%%%%%%%%%%%%%%%%%%%%%
%
% List of Tables
\newpage
\renewcommand{\listtablename}{\protect\makebox[5cm][s]{\nameLot}}
%\makebox{} is fragile; need protect
\phantomsection % for hyperref to register this
\addcontentsline{toc}{chapter}{\nameLot}
\listoftables

%%%%%%%%%%%%%%%%%%%%%%%%%%%%%%%
%       符號說明 
%%%%%%%%%%%%%%%%%%%%%%%%%%%%%%%
%
% Symbol list
% define new environment, based on standard description environment
% adapted from p.60~64, <<The LaTeX Companion>>, 1994, ISBN 0-201-54199-8

\newpage
%% 論文本體頁碼回復為阿拉伯數字計頁，並從頭起算
\pagenumbering{arabic}
%%%%%%%%%%%%%%%%%%%%%%%%%%%%%%%%